\section{An evaluated experiment and its sharp distortion constant}
\label{sec:v27-evaluated}
The separator bound retains acquisition mass, coordinate density and query recovery as separate quantities. A comparison of event masses alone therefore does not compare the resulting risk bounds. We evaluate all three quantities in an actual experiment, and also determine its optimal high-resolution average distortion. This last calculation distinguishes a useful lower certificate from an exact leading constant.

\subsection{The two-trial experiment}
Take the graph with two vertices and one edge, $n=m=1$, the uniform prior on $[0,1]$, and the positive cells
\[
                 k_0(t)=\frac12+\frac t4,\qquad
                 k_1(t)=\frac12-\frac t4.
\]
The acquisition command $(x,y)$ is uniform on $[1/4,3/4]^2$, with density $4$. Its three report likelihoods are $xk_0$, $yk_1$ and
\[
 L_{\mathrm f}(t;x,y)=1-xk_0(t)-yk_1(t)
                  =1-\frac{x+y}{2}+\frac{y-x}{4}t.
\]
After the first trial the independently selected query uses one of the four corner commands, uniformly. The second trial is executed for every query and every first report. Thus every physical run has exactly two trials and one queried boundary.

Let $z=\E[t\mid h]$. In the order $(1/4,1/4)$, $(1/4,3/4)$, $(3/4,1/4)$, $(3/4,3/4)$, the future all-failure probability vector is
\begin{equation}\label{eq:v27-query-line}
            q(z)=\left(\frac34,\frac12+\frac z8,
                                  \frac12-\frac z8,\frac14\right).
\end{equation}
With the uniform query norm $\|u\|_q^2=\tfrac14\sum_{i=1}^4u_i^2$,
\begin{equation}\label{eq:v27-query-isometry}
               \|q(z)-q(z')\|_q^2=(z-z')^2/128.
\end{equation}
The affine recovery $z=4(q_2-q_3)$ is defined on arbitrary centres; its linear operator norm from $\|\cdot\|_q$ to the real line is $L=8\sqrt2$.

\begin{proposition}[The full law of the acquired mean]
\label{prop:v27-density}
The unconditional law of $z$ is
\begin{equation}\label{eq:v27-full-law}
 \nu=\frac5{16}\delta_{8/15}+\frac3{16}\delta_{4/9}+f(z)\,dz,
\end{equation}
where the continuous part is supported on $[10/21,14/27]$. With $u=48(z-1/2)$, its density is
\begin{equation}\label{eq:v27-density}
 f(z)=
 \begin{cases}
 512\left\{\dfrac{27}{(8-5u)^3}-\dfrac1{(8+3u)^3}\right\},
                                      &-8/7\le u\le0,\\[6pt]
 512\left\{\dfrac{27}{(8+3u)^3}-\dfrac1{(8-5u)^3}\right\},
                                      &0\le u\le8/9,\\[4pt]
 0,&\text{otherwise}.
 \end{cases}
\end{equation}
It is continuous, increases from zero to $26$ on the first interval, decreases to zero on the second, and has total mass $1/2$.
\end{proposition}
\begin{proof}
On report $0$, the command factor $x$ cancels in Bayes' formula. The acquired mean is $(\int_0^1tk_0(t)\,dt)/(\int_0^1k_0(t)\,dt)=(1/3)/(5/8)=8/15$, and the report mass is $(1/2)(5/8)=5/16$. Report $1$ similarly gives mean $(1/6)/(3/8)=4/9$ and mass $3/16$.

On a failure, let
\[
             E(x,y)=1-\frac58x-\frac38y.
\]
The acquired mean is
\begin{equation}\label{eq:v27-failure-mean}
               z=\frac12+\frac{y-x}{48E(x,y)}.
\end{equation}
This is increasing in $y$ and decreasing in $x$; its extreme corner values are $10/21$ and $14/27$. The transformation from $(x,u)$ to $(x,y)$ is
\[
 y=\frac{8u+(8-5u)x}{8+3u},\qquad
 E=\frac{8(1-x)}{8+3u},\qquad
 \frac{\partial y}{\partial z}
                        =\frac{3072(1-x)}{(8+3u)^2}.
\]
The failure-weighted command density is $4E(x,y)$. Consequently its pushforward to $(x,z)$ has density
\[
                         \frac{98304(1-x)^2}{(8+3u)^3}.
\]
For $0\le u\le8/9$, the admissible $x$ interval is
\[
       \frac14\le x\le\frac{24-23u}{4(8-5u)};
\]
for $-8/7\le u\le0$, it is
\[
       \frac{8-29u}{4(8-5u)}\le x\le\frac34.
\]
Integrating the quadratic numerator over these intervals gives \eqref{eq:v27-density}. The endpoints give zero and both middle limits give $26$. Differentiation with respect to $u$ gives, on the first interval,
\[
 512\left\{\frac{405}{(8-5u)^4}+\frac9{(8+3u)^4}\right\}>0,
\]
and on the second interval
\[
 -512\left\{\frac{243}{(8+3u)^4}+\frac{15}{(8-5u)^4}\right\}<0.
\]
Finally, integrating the actual failure evidence over the command square gives $1-(5/8)(1/2)-(3/8)(1/2)=1/2$. No report evidence has been removed from this density.
\end{proof}

\subsection{Exact reduction and the leading minimax constant}
For a finite nonnegative measure $\lambda$ on the line define
\[
 D_M(\lambda)=\inf_{\substack{C\subset\R\\1\le|C|\le M}}
                    \int\operatorname{dist}(z,C)^2\,d\lambda(z).
\]
This definition does not normalize the mass of $\lambda$.

\begin{proposition}[Exact scalar quantization reduction]
\label{prop:v27-exact-reduction}
For every integer $M\ge1$, the optimal average excess in this graph experiment, including randomized controllers, is
\begin{equation}\label{eq:v27-exact-reduction}
                   \mathcal R^{\mathrm{av}}_{G,M}=D_M(\nu)/128.
\end{equation}
The same equality holds in the full-history scheduler and finite-prefix decoder relaxation. The quantity in \eqref{eq:v27-exact-reduction} is an average-risk optimum; no equality with the maximum-tape optimum is asserted.
\end{proposition}
\begin{proof}
Project any prediction centre orthogonally, in the query norm, onto the affine line \eqref{eq:v27-query-line}. Its distance to every true query vector cannot increase. Clipping the corresponding scalar centre to the convex hull of the support of $\nu$ further decreases distance and yields admitted probability predictions. For any at most $M$ centres, a randomized assignment cannot beat the nearest-centre assignment. By \eqref{eq:v27-query-isometry}, the infimum of their average losses is $D_M(\nu)/128$.

Conversely, the single acquired batch determines $z$ exactly for purposes of the transition computation. An $M$-label transition can assign it to a nearest scalar centre, retain only that index, and output \eqref{eq:v27-query-line} at the query. Thus the scalar infimum is attained or approached by admissible finite-state predictors; temporary computation is not the retained resource.

The first vertex is selected before any command or report is seen. Its order and any independent seed therefore provide no tape information. Conditional on them, the same scalar law and the same lower bound apply. Averaging cannot lower the infimum. The two deterministic orders have the same experiment, and there is only one boundary, so both the original and oracle risks have the stated value.
\end{proof}

We include the one-dimensional high-resolution argument to make the leading-constant calculation independent of any unstated regularity theorem. Its general quantization content is classical; see \cite{GrafLuschgy2000}. Its application below uses the actual mixed acquired law, rather than an assumed uniform latent statistic.

\begin{lemma}[Continuous density with finitely many atoms]
\label{lem:v27-quantization}
Let $g$ be continuous and nonnegative on a compact interval and zero off that interval. If $\lambda=g(z)\,dz+\sum_{r=1}^J a_r\delta_{z_r}$, with $J<\infty$ and $a_r\ge0$, then
\begin{equation}\label{eq:v27-quantization}
          \lim_{M\to\infty}M^2D_M(\lambda)
                   =\frac1{12}\left(\int g(z)^{1/3}\,dz\right)^3.
\end{equation}
\end{lemma}
\begin{proof}
First consider $g(z)\,dz$. Partition its supporting interval into $L$ subintervals $I_i$, of lengths $l_i$, and put $h_i=\inf_{I_i}g$. The Voronoi intervals of an at most $M$-point codebook intersect this partition in at most $M+L-1$ nonempty intervals. If $n_i$ intersections occur in $I_i$, their lengths sum to $l_i$. On an interval of length $d$, the integral of squared distance to any fixed centre is at least $d^3/12$. Convexity of the cube and then H\"older's inequality give
\[
 \begin{split}
 \int\operatorname{dist}(z,C)^2g(z)\,dz
 &\ge\frac1{12}\sum_i\frac{h_i l_i^3}{n_i^2}\\
 &\ge\frac{(\sum_i h_i^{1/3}l_i)^3}
                         {12(M+L-1)^2}.
 \end{split}
\]
This estimate holds for every codebook. Let $M$ tend to infinity at fixed partition, then refine the partition. Continuity of $g^{1/3}$ proves the lower bound in \eqref{eq:v27-quantization}.

For the upper bound put $H_i=\sup_{I_i}g$. At a fixed partition assign $m_i$ equally spaced midpoint centres to $I_i$, with total at most $M$ and with $m_i/M$ tending to $H_i^{1/3}l_i/\sum_r H_r^{1/3}l_r$ for every positive summand. Such integer allocations exist by rounding and reserving at most $L$ centres. Their assigned-cell error is at most
\[
                         \sum_i\frac{H_i l_i^3}{12m_i^2}.
\]
Nearest-centre assignment among all the centres is no worse. Taking the limit at this fixed partition and then refining gives the upper bound. Intervals with $H_i=0$ contribute no error; the case $g=0$ is immediate.

Adding nonnegative atoms cannot improve the lower bound for the continuous submeasure. For an upper bound place a centre at each of the finitely many atoms and use the remaining $M-J$ centres for the continuous part. Since $M/(M-J)\to1$, the same limit follows. This also covers atoms outside the continuous support.
\end{proof}

\begin{theorem}[Sharp average distortion for the evaluated experiment]
\label{thm:v27-sharp-constant}
For the positive two-trial experiment above,
\begin{equation}\label{eq:v27-sharp-constant}
 \lim_{M\to\infty}M^2\mathcal R^{\mathrm{av}}_{G,M}
     =\kappa_{\mathrm{ex}}:=
       \frac1{1536}\left(\int_{10/21}^{14/27}f(z)^{1/3}\,dz\right)^3.
\end{equation}
Numerically $\kappa_{\mathrm{ex}}$ is approximately $4.886\,10^{-7}$. The integral expression, not the decimal approximation, is the exact constant.
\end{theorem}
\begin{proof}
Combine Propositions~\ref{prop:v27-density} and \ref{prop:v27-exact-reduction} with Lemma~\ref{lem:v27-quantization}. The factor is $1/(128\cdot12)=1/1536$.

For reproducible bounds on the numerical value, split the $u$ interval at zero. The function $f^{1/3}$ is monotone on each part. Uniform lower and upper endpoint sums, with $dz=du/48$, bound its integral. Every endpoint density is rational. For a rational $r\ge0$ and an integer scale $s_0$, compute the integer cube root of $\lfloor r s_0^3\rfloor$; if it is $b$, then $b/s_0\le r^{1/3}<(b+1)/s_0$. Thus both endpoint sums admit entirely rational bounds. Cubing them and dividing by $1536$ bounds \eqref{eq:v27-sharp-constant} without treating floating-point quadrature as a proof. The supplied diagnostic implements this construction with $8192$ intervals on each side and scale $10^{12}$.
\end{proof}

\subsection{A complete same-model separator comparison}
Restrict the first command to
\[
          \Omega=[1/4,3/8]\times[5/8,3/4]
\]
and require its failure report. Denote this first-block event by $A$. Direct integration over the commands gives the latent likelihood subprobability polynomial
\[
                  \frac1{32}+\frac{3t}{512}.
\]
In particular,
\begin{equation}\label{eq:v27-beta}
 \beta=\Prob(A)=\frac{35}{1024},\qquad
 \E[z\mid A]=\frac{18}{35}.
\end{equation}
The latter equality is an average over this actual subevent, not a fixed value of the acquired mean.

\begin{proposition}[Evaluated density and recovery costs]
\label{prop:v27-certificate}
The pushforward of the first-block subprobability $\Prob(A,\cdot)$ to $z$ has density at most
\begin{equation}\label{eq:v27-H}
 H=\frac{192}{5}\left(\frac{39}{64}\right)^3
                         =\frac{177957}{20480}.
\end{equation}
With $L=8\sqrt2$ from \eqref{eq:v27-query-line} and
\[
                 C=\frac{2LH}{\beta},\qquad
                 C^2=\frac{1013398203168}{30625},
\]
the full separator certificate is, for every $M\ge1$,
\begin{equation}\label{eq:v27-evaluated-certificate}
 \mathcal R^{\mathrm{av}}_{G,M}\ge
       \frac{\beta}{12C^2M^2}
       =\frac{1071875}{12452637120528384}\,M^{-2}.
\end{equation}
If instead one also requires the completion report to be a failure, while retaining exactly the same decoder, query norm and recovery map, the corresponding complete certificate is exactly one half of \eqref{eq:v27-evaluated-certificate}.
\end{proposition}
\begin{proof}
On $\Omega$, the first-word evidence obeys $31/64\le E(x,y)\le39/64$. Differentiating \eqref{eq:v27-failure-mean} gives
\[
                   \frac{\partial z}{\partial y}
                               =\frac{1-x}{48E(x,y)^2}>0.
\]
The failure-weighted joint density after the change of variable $y\mapsto z$ is therefore $192E(x,y)^3/(1-x)$. At a fixed $z$, its $x$ integration domain is a subset of an interval of length $1/8$, and $1-x\ge5/8$. Bounding the integrand by its stated upper values proves \eqref{eq:v27-H}. This argument evaluates a weighted subprobability density directly; it does not assume that the chosen rectangle is a particular inverse chart from Lemma~\ref{lem:v25-regular-box}.

At either possible internal visited set the decoder has at most $M$ centres. By the recovery norm, loss at most $t$ puts $z$ in a union of at most $M$ intervals of radius $L\sqrt t$. Intersecting with the visited-set event only decreases the unconditioned mass. Thus its node capacity is bounded by
\[
                   c(t)=\min\{1,MC\sqrt t\}.
\]
Both internal vertices are needed to separate the source and sink of the two-vertex subset lattice, so $H_G(t)=2c(t)$. The separator proof applies to this directly evaluated density just as it does to the chart density. Since there is one queried boundary,
\[
 \begin{split}
 \mathcal R^{\mathrm{av}}_{G,M}
 &\ge\beta\int_0^\infty[1-2c(t)]_+\,dt\\
 &=\beta\int_0^{1/(4M^2C^2)}(1-2MC\sqrt t)\,dt
   =\frac{\beta}{12M^2C^2}.
 \end{split}
\]
Substitution of the rational values gives \eqref{eq:v27-evaluated-certificate}.

For the comparison event, integrate the completion command over the same uniform square. Its failure likelihood conditional on any latent $t$ is
\[
                  \E_{x,y}[1-xk_0(t)-yk_1(t)]=\frac12.
\]
Consequently requiring that completion failure multiplies the entire first-block subprobability measure, not merely its total mass, by $1/2$. Its mass is $\beta/2$ and a valid density bound is $H/2$. The conditional density ratio $H/\beta$, the recovery norm $L$, and hence every capacity $c(t)$ are unchanged. The prefactor in the integrated certificate is halved. This is a like-for-like comparison with the same decoder; it does not transfer constants between the prefix-oracle and original decoder models.
\end{proof}

The coefficient in \eqref{eq:v27-evaluated-certificate} is approximately $8.608\,10^{-11}$, whereas the sharp average coefficient in \eqref{eq:v27-sharp-constant} is approximately $4.886\,10^{-7}$. Their ratio is about $5.68\,10^3$. The separator certificate is therefore conservative in this instance, even though removing the completion-word restriction improves the complete certificate by a rigorously evaluated factor of two. The sharp-constant result is specific to this experiment; it is not a claim that the general separator bound supplies sharp constants for all graphs or collision configurations.
